\documentclass[a4paper, twocolumn]{article}
\usepackage{graphicx} % Required for inserting images
\usepackage{lipsum}
\usepackage{hyperref}
\usepackage{amsmath}
\usepackage[style=science]{biblatex}
\addbibresource{refs.bib}

\title{Tutorial Exercise 2: References and Citations with \\  Bib\LaTeX\ (Submission Document)}
\author{K Kosgahakumbura}
\date{September 2026}

\begin{document}
\maketitle

\section{Introduction}
This document is a sample solution for Tutorial Exercise 2. Its brief is available as a PDF on our Moodle page

\section{Literature Research}

When writing your literature review and final
project dissertation, base your work mostly on recent peer-reviewed references, as these are reliable
sources.

For example, when writing a literature review
on Evolutionary Algorithms (EA) applied to Machine Learning (ML), reading the Wikipedia page
on EA \cite{wikipedia} or a blog post \cite{adan}
could help you get
started with the basics and provide pointers to
peer-reviewed references. However, you should not
cite such web pages because they lack editorial
oversight and quality control, as non-experts may
contribute content and cannot be verified for accuracy. In other words, these are not peer-reviewed
sources

Textbooks \cite{de2016evolutionary} can provide a good foundation in
knowledge and context, but be cautious about citing older editions, as the field evolves rapidly. Additionally, not all textbooks undergo the same rigorous peer-review process as journal articles. Always
prioritise citing recent peer-reviewed papers,
particularly primary sources that report original research.

Secondary sources, such as survey \cite{akbar} or review
papers, are valuable for understanding the broader
context, as they synthesise existing research to help
identify key themes and primary sources. How
ever, secondary research citations should complement,
not replace, primary ones. For example,
while \textcite{akbar} investigate the role of EA in
solving different ML challenges, \textcite{zhang}
propose a new generator of optimisation benchmark
problems for EA, and define a fitness function as

\begin{equation}
  F = \alpha A +\beta B + \gamma P
    \label{eq:equation}
\end{equation}
where $\alpha$,  $\beta$, and $\gamma$  are
weight coefficients. The term
A in (\ref{eq:equation}) quantifies the discrepancy in the global
optimum position of the generated and stipulated
landscapes, while B measures the positional error
between the peaks of these landscapes. Mathematically,
$P = \sqrt{ (2L)2 + (2U )2}$
denotes the maximum Euclidean distance within the landscape.

ArXiv preprints can provide access to recent developments
\cite{alexander}, but should be used cautiously and
kept to a minimum, as they’re not peer-reviewed.
Always prioritise peer-reviewed sources as the foundation of your literature review. ArXiv papers
may be appropriate only when they represent genuinely
influential work with no peer-reviewed equivalent, but ensure they don’t dominate your references.

\printbibliography
\end{document}

